% DRAFT (plan 14.4, 14.7). Not loaded by main.tex. It goes into the book at the book pass, once the five-laboratory test (bar 2) and
% layer 1 have run. Every number below is printed by the chain in a committed file:
%   doors/data/DEV_METAA_HEALTHY_01/synthetic_02_summary.txt (run_sample.py on a real healthy blood and on synthetic arrays)
%   doors/data/DEV_METAA_HEALTHY_01/gain_counter_01_summary.txt (Stage Q2's gain-side count on constructed molecules)
% Still to write before it goes in: bar 2 (100 healthy bloods, five laboratories), layer 1 (the 56 cell types), the per-cell gap, figures.

\section{Where Met-A's healthy middle line comes from}\label{sec:metaa_healthy_line}

IAM-A's middle line is derived: a sample's own loose share predicts the error a healthy cell of its kind should make, and the reading is
the measured error over that prediction. Met-A's middle line has so far been measured: the disorder a healthy sample of the same kind
reads on the same platform, or IAM-Atlas v2.1 mixed in the sample's own cell shares. This section is the route from measuring that line
to predicting it, in two layers.

\subsection{Why a healthy blood does not read 1}

The same healthy blood reads differently in different laboratories. Before any correction, healthy whole blood reads about 1.11 in one
series and about 0.96 in another. The blood is not the difference; the arrays are. An array does not report a site's true level. It
reports a compressed, noisy copy: a probe design's lowest readings sit a little above 0, its highest a little below 1, and every reading
carries noise. Disorder is curved, and it rises steeply near 0 and near 1, so even a small offset at a site that is almost never marked
adds disorder that is not in the blood.

\subsection{Layer one: the array measures its own blur}

Some sites are held the same way in all 79 healthy cell types of the atlas, either firmly marked or firmly unmarked. Whatever an array
reads at those anchor sites is the array, not the blood. On them the chain measures, for each probe design, the array's own straight line
(where its 0 and its 1 actually fall) and its own noise. From those three numbers it predicts what the array would read if the sample
were exactly healthy, the atlas mixed in the sample's own shares, and divides the measured reading by that prediction. Nothing is fitted:
the line and the noise are measurements of the array where every healthy cell agrees.

On one real healthy blood read from its raw files, the prediction takes the reading from 1.11291 to 1.0127. Through the whole chain, on
synthetic arrays built with a known blur, a healthy sample reads 0.99689, 0.99866 and 0.99278 on mild, typical and weak arrays, and a loss
placed away from the anchors keeps most of its rise. One limit is physical and stays: a loss spread evenly over every site, the anchors
included, looks exactly like a weaker array, and no reading made on the array alone can tell the two apart. Cancer's losses are not even;
they fall in the loosely held stretches of the genome, not at sites every cell type holds firmly.

\subsection{Layer two: the atlas's healthy levels from the held error}

The prediction still uses the atlas's measured healthy levels. The next layer asks where those levels come from. A cell that holds a mark
at a site keeps it on every copy except where a copy drops it, so its level there should be one minus its held error on the loss side. A
cell that holds a site unmarked should sit at its held error on the gain side, how often a copy adds a mark it should not. Stage Q2 now
counts both on a cell's own molecules; on constructed molecules it recovers known rates within a few percent, the same on both sides.
Whatever is left after the prediction is the gap: sites a cell holds only partly, the part only biology sets, measured once. If the
prediction holds, the atlas need only say which sites each cell holds; the levels come from physics and the blur from the array.
